%% Documento tecnico de planejamento. Registro interno (PAVIC/UFPI).
\documentclass[12pt,a4paper]{article}

\usepackage[utf8]{inputenc}
\usepackage[T1]{fontenc}
\usepackage[brazil]{babel}
\usepackage[margin=2.4cm]{geometry}
\usepackage{booktabs}
\usepackage{enumitem}
\usepackage{titlesec}
\usepackage{graphicx}
\usepackage{subcaption}
\graphicspath{{figs_dados/}}

\setlength{\parindent}{0pt}
\setlength{\parskip}{0.55em}
\titlespacing*{\section}{0pt}{1.5em}{0.4em}
\titlespacing*{\subsection}{0pt}{0.95em}{0.25em}
\renewcommand{\arraystretch}{1.3}

\begin{document}

{\LARGE\bfseries Avaliação de carcaças ovinas por imagem\par}
\vspace{0.4em}
{\small PAVIC/UFPI}

\vspace{0.6em}
\hrule
\vspace{0.8em}

\textbf{Resumo.} Avaliaram-se mais de trinta abordagens computacionais, em quatorze
frentes experimentais, para estimar a tipificação de carcaças ovinas a partir de
imagens. Nenhuma abordagem predisse o grau a partir da imagem. O resultado central deste
registro é que o progresso foi impedido pela integridade dos dados disponíveis, e não
por insuficiência dos métodos. O pareamento entre imagem, animal e grau não foi
registrado de forma verificável, e as associações foto--nota utilizadas são incertas. Em
consequência, toda correlação entre descritores de imagem e o grau é não interpretável,
o que invalida os resultados de predição e impede testar qualquer hipótese sobre o que a
imagem contém. As únicas saídas que não dependem do grau, a segmentação de tecidos e a
consistência anatômica do mapa de convexidade do perfil, permaneceram disponíveis e
operaram conforme o esperado. Este documento descreve a causa da inviabilidade dos
dados, os métodos avaliados com os respectivos resultados, e os requisitos de uma nova
aquisição. Os resultados são limitados pela amostra disponível (n = 22 carcaças ovinas,
anotação manual de tecidos) e pela falha de pareamento descrita.

\tableofcontents
\vspace{0.6em}
\hrule

%========================================================================
\section{Por que os dados disponíveis são inviáveis}

A inviabilidade tem uma causa demonstrada: a falha de integridade do pareamento entre
imagem, animal e grau. Essa falha, isoladamente, impede a estimação do grau a partir da
imagem. Ela também impede testar qualquer hipótese sobre o conteúdo da imagem, pois todo
teste depende de um grau corretamente pareado.

\subsection{Falha de integridade do pareamento imagem--animal--grau}

O pareamento entre cada fotografia e o respectivo animal e grau não foi registrado no
momento da coleta. Ele foi reconstruído posteriormente por ordenação dos nomes de
arquivo e associação sequencial à lista de animais, sem verificação. Três evidências,
extraídas dos dados, confirmam a inconsistência.

\textbf{Evidência 1: etiquetas físicas com duplicatas e ilegibilidade.} As etiquetas
físicas das carcaças foram lidas manualmente sobre as imagens. A leitura indica que três
imagens consecutivas correspondem à mesma carcaça, e que uma etiqueta é ilegível
(Tabela~\ref{tab:tags}, Figura~\ref{fig:dup}). A numeração das carcaças nas etiquetas,
além disso, não segue a ordem dos arquivos de imagem.

\begin{figure}[!ht]
\centering
\begin{subfigure}{0.235\textwidth}\centering
\includegraphics[width=\textwidth]{IMG_6194_small.jpg}
\caption{IMG\_6194}
\end{subfigure}\hfill
\begin{subfigure}{0.235\textwidth}\centering
\includegraphics[width=\textwidth]{IMG_6195_small.jpg}
\caption{IMG\_6195}
\end{subfigure}\hfill
\begin{subfigure}{0.235\textwidth}\centering
\includegraphics[width=\textwidth]{IMG_6196_small.jpg}
\caption{IMG\_6196}
\end{subfigure}\hfill
\begin{subfigure}{0.235\textwidth}\centering
\includegraphics[width=\textwidth]{IMG_6203_small.jpg}
\caption{IMG\_6203}
\end{subfigure}
\caption{Exemplos das imagens da coleta. (a)--(c) IMG\_6194, IMG\_6195 e IMG\_6196 são a
mesma carcaça (etiqueta legível, carcaça 10), registrada três vezes; um pareamento
sequencial as associaria a três animais distintos. (d) IMG\_6203 tem a etiqueta (saco
plástico no jarrete) em branco, ilegível. As imagens também ilustram a ausência de
padronização de captura: fundo não uniforme, presença do operador, e iluminação por
janela.}
\label{fig:dup}
\end{figure}

\begin{table}[!ht]
\centering
\caption{Trecho da leitura manual das etiquetas físicas. A coluna ``carcaça'' é o
identificador escrito na etiqueta; a coluna ``tratamento'' é a dieta registrada na
etiqueta. As três imagens IMG\_6194 a IMG\_6196 correspondem à carcaça 10, e a IMG\_6203
tem etiqueta ilegível.}
\label{tab:tags}
\begin{tabular}{lll}
\toprule
Imagem & Carcaça (etiqueta) & Tratamento (etiqueta) \\
\midrule
IMG\_6193 & 01 & T1 / 0\% \\
IMG\_6194 & 10 & T1 / 0\% \\
IMG\_6195 & 10 & T1 / 0\% \\
IMG\_6196 & 10 & T1 / 0\% \\
IMG\_6203 & ilegível & ilegível \\
IMG\_6204 & 14 & T1 / 0\% \\
IMG\_6217 & 28 & T1 / 0\% \\
\bottomrule
\end{tabular}
\end{table}

\textbf{Evidência 2: a numeração da etiqueta não corresponde à do laboratório.} O
identificador da etiqueta da carcaça e o identificador (Animal) da planilha de
laboratório constituem sistemas de numeração distintos, sem tabela de correspondência.
A verificação compara, por número, o tratamento da etiqueta com o tratamento que o
laboratório atribui ao Animal de mesmo número (Tabela~\ref{tab:cruz}). Das 15 carcaças
de etiqueta legível e presentes no laboratório, apenas 4 coincidem, no nível do acaso.
As carcaças 24 e 28 não constam da planilha. A carcaça 16 tem tratamento T5 (30\% mais
TAN), ausente da planilha. Sete animais da planilha (5, 9, 11, 13, 18, 19, 22) não
possuem imagem correspondente.

\begin{table}[!ht]
\centering
\caption{Verificação cruzada do tratamento por número. ``\% (etiqueta)'' é a dieta lida
na etiqueta da carcaça; ``\% (laboratório)'' é a dieta que a planilha atribui ao Animal
de mesmo número. Coincidem 4 de 15 casos testáveis.}
\label{tab:cruz}
\begin{tabular}{cccc}
\toprule
Número & \% (etiqueta) & \% (laboratório) & Coincide \\
\midrule
01 & 0  & 0  & sim \\
10 & 0  & 10 & não \\
14 & 0  & 30 & não \\
03 & 10 & 10 & sim \\
06 & 10 & 10 & sim \\
08 & 10 & 20 & não \\
15 & 10 & 20 & não \\
04 & 20 & 30 & não \\
17 & 20 & 30 & não \\
23 & 20 & 10 & não \\
02 & 30 & 20 & não \\
07 & 30 & 0  & não \\
12 & 30 & 20 & não \\
21 & 30 & 30 & sim \\
16 & 30+TAN & 0 & não \\
\bottomrule
\end{tabular}
\end{table}

\textbf{Evidência 3: a planilha de laboratório não referencia as imagens.} A planilha de
laboratório possui 22 linhas e 27 colunas. A coluna de identificação é o Animal, e não
há coluna de nome de imagem. Não existe, portanto, vínculo registrado entre a planilha e
os arquivos de foto.

Como consequência, toda correlação entre descritores de imagem e o grau utilizou
rótulos de associação incerta. O mesmo se aplica às correlações com a espessura de
gordura e com a área de olho de lombo. Esses resultados não são interpretáveis quanto
ao grau.

\subsection{Hipótese de modalidade, não testável com os dados atuais}

A tipificação considera dois eixos. A conformação descreve a forma e o desenvolvimento
muscular, e é definida sobre o perfil externo. O acabamento descreve a cobertura de
gordura subcutânea, e depende em parte da espessura da gordura.

A distinção entre os dois eixos motiva uma hipótese sobre o alcance da imagem. A
conformação é uma propriedade do perfil externo, em princípio registrada na projeção de
superfície. A espessura da gordura é uma propriedade de profundidade, que por
construção óptica não é registrada na projeção de superfície (Tabela~\ref{tab:foto}).
Esta hipótese é consistente com a prática reportada na literatura, em que a composição
de gordura é estimada por modalidades de profundidade, como a absorciometria de raios X,
a tomografia e o ultrassom, e não por imagem de superfície (Seção~\ref{sec:lit}).

Cabe a ressalva central: esta hipótese não foi testada neste projeto. Os testes de
acabamento dependeram do grau de pareamento incerto, e portanto não constituem
evidência a favor nem contra a hipótese. A separação entre o que é uma expectativa
óptica e o que seria um resultado próprio do projeto requer uma aquisição com pareamento
verificável.

\begin{table}[!ht]
\centering
\caption{Atributos observáveis e não observáveis na imagem de superfície da carcaça.}
\label{tab:foto}
\begin{tabular}{p{6.6cm}p{6.6cm}}
\toprule
Observável na imagem & Não observável na imagem \\
\midrule
perfil da perna, do lombo e da paleta (convexo ou côncavo) & espessura da gordura de
cobertura \\
largura e dimensões da carcaça & gordura perirrenal \\
exposição óssea no contorno & marmoreio (gordura intramuscular) \\
\bottomrule
\end{tabular}
\end{table}

%========================================================================
\section{Métodos avaliados e resultados}

Avaliaram-se mais de trinta abordagens. Para a exposição, agrupam-se em cinco categorias
segundo o resultado obtido. Os coeficientes reportados situam-se no intervalo $[0,1]$,
em que 1 indica associação perfeita e 0 indica ausência de associação; n = 22 salvo
indicação em contrário.

\subsection{Segmentação de tecidos (etapa de pré-processamento)}

A segmentação separa, na imagem, as regiões de gordura, músculo e fundo. Ela antecede
qualquer medição. Avaliaram-se uma tabela de correspondência cromática (ColorNaming),
uma rede pré-treinada (DINOv2) e a replicação de uma arquitetura da literatura. As
regiões de fundo e gordura foram segmentadas com sobreposição (IoU) próxima de 0,83 para
a gordura. A região de músculo apresentou sobreposição inferior, de 0,10 a 0,38. Essa
diferença é consistente com a similaridade cromática entre músculo e fundo. A etapa não
depende do grau e atende ao requisito de pré-processamento.

\subsection{Estimação do acabamento (não interpretável sob os dados atuais)}

Avaliaram-se descritores de área e de fração de gordura da máscara, estimadores de
espessura por transformada de distância, e descritores cromáticos. Avaliaram-se também
métodos geométrico-estatísticos: homologia persistente (conectividade e número de
cavidades da cobertura de gordura), análise de frequência e textura, percolação, e
descritores de simetria e fragmentação. A maior correlação de Spearman com o acabamento
foi $\rho \approx 0{,}49$, obtida com o descritor de área, e nenhum descritor atingiu
$|\rho| \geq 0{,}5$.

Esses valores não são interpretáveis, pois foram calculados contra o grau de pareamento
incerto descrito na Seção~1. Um valor baixo, sob rótulos incertos, é compatível tanto
com uma ausência real de associação quanto com o efeito do pareamento. Por esse motivo,
os resultados de acabamento não constituem evidência sobre o alcance da imagem. A
hipótese de modalidade da Seção~1.2 permanece não testada.

\subsection{Estimação da conformação (resultados não conclusivos sob os dados atuais)}

Avaliaram-se descritores de morfometria de silhueta (largura, extensão, solidez),
descritores elípticos de Fourier do contorno, e o invariante integral de convexidade do
perfil. Diversos descritores apresentaram correlações brutas de 0,5 a 0,6. Essas
correlações não se sustentaram após a correção para comparações múltiplas. Parte delas
foi atribuída a uma fonte de imagem incorreta, o contorno da região de corte em lugar do
contorno da carcaça. As demais utilizaram o pareamento incerto. Um teste de limite
superior foi conduzido sobre o contorno anotado manualmente. Nenhum dos 384 testes
permaneceu significativo após a correção de Benjamini--Hochberg.

Esse resultado requer interpretação cautelosa, pois foi obtido contra o mesmo grau de
pareamento incerto. Ele não estabelece que a conformação seja inobservável na imagem.
Indica, antes, que a hipótese não é testável sob os dados atuais. A conformação permanece
a frente compatível com a modalidade bidimensional.

\subsection{Métodos sem desempenho acima da referência}

Alguns métodos não superaram a referência do acaso, independentemente da integridade dos
dados. Esse desempenho é consistente com o tamanho de amostra reduzido (n = 22). O grupo
inclui a transformada de espalhamento por \emph{wavelets}, com acurácia balanceada de
0,14 ante uma referência de 0,33, a transformada de característica de Euler, a análise de
renormalização por suavização em escala, e a seleção esparsa de variáveis.

\subsection{Resultado independente do grau}

O mapa de convexidade do perfil é anatomicamente consistente: as pernas e os jarretes
são classificados como convexos, e o entalhe medial entre as pernas como côncavo. Este
resultado não depende do grau e indica que a forma é mensurável na imagem. A validação
preditiva requer um grau corretamente pareado, ainda indisponível.

\subsection{Evidência externa da literatura}\label{sec:lit}

Uma revisão de aproximadamente 150 fontes fornece evidência externa, sobre dados de
pareamento correto, distinta dos resultados não interpretáveis acima. Os sistemas de
análise de imagem (VIA) reportam correlação de 0,94 a 0,97 para conformação e valores
inferiores e mais dispersos (0,30 a 0,91) para gordura. A composição de gordura é
estimada por absorciometria de raios X (DEXA), tomografia ou ultrassom. Essa evidência
é consistente com a hipótese de modalidade da Seção~1.2. Ela não substitui um teste
sobre os dados próprios do projeto, ainda pendente. A Tabela~\ref{tab:resumo} resume as
cinco categorias.

\begin{table}[!ht]
\centering
\caption{Síntese dos métodos avaliados, por categoria de resultado.}
\label{tab:resumo}
\begin{tabular}{p{3.5cm}p{4.5cm}p{4.8cm}}
\toprule
Categoria & Resultado & Causa \\
\midrule
Segmentação de tecidos & gordura e fundo segmentados; músculo inferior & etapa de
pré-processamento; atende \\
Acabamento & $\rho \leq 0{,}49$ (não interpretável) & grau de pareamento incerto \\
Conformação & correlações não sustentadas após correção & grau de pareamento incerto \\
Métodos geométrico-estatísticos & no nível da referência ou abaixo & amostra reduzida
(n = 22) \\
Mapa de convexidade & anatomicamente consistente & não depende do grau; requer grau
pareado para validar \\
\bottomrule
\end{tabular}
\end{table}

%========================================================================
\section{Requisitos para a continuidade}

A reanálise dos dados atuais não é indicada para a estimação do grau. Requer-se uma nova
aquisição com quatro requisitos.

\begin{enumerate}[itemsep=0.35em,topsep=0.3em]
\item \textbf{Pareamento físico entre imagem, animal e medidas, no momento do abate.} A
etiqueta da carcaça deve ser legível na própria imagem, e uma única planilha,
preenchida na coleta, deve associar a cada carcaça as suas imagens, notas e medidas. O
pareamento não deve ser reconstruído posteriormente.

\item \textbf{Grau atribuído por avaliadores sobre a própria imagem.} Dois a três
avaliadores treinados atribuem a nota de forma independente, a partir da mesma imagem
utilizada na análise. Mede-se a concordância entre avaliadores e deriva-se um grau de
consenso. A concordância insuficiente invalida o uso da nota como referência.

\item \textbf{Amostra suficiente e estratificada.} Estima-se um total de 100 a 120
carcaças, cobrindo a faixa completa de grau, das categorias de menor conformação e
acabamento às de maior. A amostra atual (n = 22), concentrada em poucas categorias, não
permite estimação com precisão adequada.

\item \textbf{Referência física, para a estimação de gordura no nível da literatura.} A
espessura de gordura, o GR, a área de olho de lombo, ou a dissecção devem ser medidos
nos mesmos animais. Sem referência física, a estimação restringe-se à conformação,
observável na imagem.
\end{enumerate}

\subsection{Limitações}

As limitações deste registro são: (1) amostra de n = 22 carcaças ovinas com anotação
manual de tecidos; (2) pareamento imagem--grau de integridade não verificada, que
invalida os resultados de predição de grau; (3) distribuição de grau concentrada em
poucas categorias; (4) ausência de validação contra referência física nos mesmos
animais; (5) resultados específicos para carcaças ovinas nas condições avaliadas.

\vspace{0.4em}
\hrule
\vspace{0.6em}

\textbf{Conclusão.} O conjunto de métodos avaliados não estabelece a predição do grau a
partir da imagem. A causa demonstrada é a integridade do pareamento entre imagem, animal
e grau. Sob rótulos de associação incerta, os resultados de predição não são
interpretáveis, e nenhuma hipótese sobre o conteúdo da imagem pode ser testada. O
progresso foi, portanto, impedido pela qualidade dos dados disponibilizados, e não por
insuficiência dos métodos. A limitação de modalidade para a gordura permanece uma
hipótese, apoiada na óptica e na literatura, mas não testada com os dados próprios. A
conformação é observável na imagem e constitui a frente compatível com a modalidade
bidimensional. A continuidade requer uma aquisição com pareamento físico verificável.

\end{document}
